\documentclass[11pt,a4paper]{article}
\usepackage[utf8]{inputenc} % Not needed with XeLaTeX
\usepackage{fontspec}
\usepackage[margin=0.75in]{geometry}
\usepackage{xcolor}
\definecolor{darknavy}{RGB}{0,0,100}
\usepackage[colorlinks=true, urlcolor=darknavy, linkcolor=darknavy]{hyperref}
\usepackage{enumitem}
% \usepackage{fontawesome}
\usepackage{fontawesome5}
\usepackage{titlesec}
\usepackage{academicons}

% % --- Font: Source Sans Pro via Type1 (works with pdfLaTeX)
% \usepackage[default,scale=1.0]{sourcesanspro} % body + sans by default
% \renewcommand{\familydefault}{\sfdefault}      % make sans the default family
% \usepackage[scaled]{helvet}
\usepackage[default,scale=0.95]{opensans}

% (optional) math that pairs reasonably with sans body
\usepackage{newtxsf} % sans math; or use newtxtext/newtxmath for serif math

\titleformat{\section}{\Large\bfseries}{\thesection}{1em}{}[\titlerule]
\titlespacing{\section}{0pt}{12pt}{6pt}

\usepackage{comicneue}

\begin{document}

\begin{center}
  % {\Huge\textbf{Pavel Popov}}\\[0.5em]
  % {\comicneue\Huge\textbf{Pavel Popov}}\\[0.5em]
  {\fontspec{JetBrains Mono}\Huge Pavel Popov}\\[0.5em]
  {\fontspec{JetBrains Mono}\footnotesize +1 (470) 920-1876 \quad | \quad Atlanta, GA 30318} \quad | \quad {paavali.popov@gmail.com}\\
  {\small \href{https://github.com/paavalipopov}{Github: paavalipopov} \quad | \quad
  % \href{https://www.linkedin.com/in/pavel-popov-83b69a123/}{\faLinkedin\ Pavel Popov} \quad
  \href{https://www.linkedin.com/in/pavel-popov-83b69a123/}{LinkedIn: pavel-popov-83b69a123} \quad | \quad
  \href{https://scholar.google.com/citations?user=4alUD5UAAAAJ}{Google scholar: 4alUD5UAAAAJ}}
\end{center}

\section*{Skills}
\begin{itemize}
[leftmargin=*,noitemsep]
\item \textbf{Data Science \& Analysis:} Exploratory data analysis (EDA), data cleaning \& preprocessing, statistical analysis, feature engineering, data visualization (Matplotlib, Seaborn, Plotly).
\item \textbf{Machine Learning:} Classification, regression, clustering, deep learning (PyTorch), model evaluation \& benchmarking, time series analysis.
\item \textbf{Programming \& Data Tools:} Python (Pandas, NumPy, SciPy, Scikit-learn), SQL, mongoDB, Git, Docker, Linux/HPC.
\item \textbf{Communication \& Collaboration:} Storytelling with data, technical writing \& documentation, cross-functional teamwork, presenting to technical \& general audiences.
\end{itemize}

\section*{Open Source Projects}
\noindent\textbf{\href{https://pypi.org/project/ml4fmri/}{ml4fmri}: Python toolkit for training \& benchmarking ML models on large-scale data}\\
[-2em]
\begin{itemize}[leftmargin=*,noitemsep]
    \item Unified 13 ML models under a standardized training/evaluation API for reproducible model comparison.
    \item Automated an end-to-end data pipeline (ingestion, preprocessing, training, evaluation) for reproducible experiments.
\end{itemize}

\noindent\textbf{\href{https://pypi.org/project/brainbow/}{brainbow} --- Data visualization tool for multi-dimensional datasets}\\
[-2em]
\begin{itemize}[leftmargin=*,noitemsep]
    \item Built a command-line tool for visualizing and exploring complex 3D/4D (neuroimaging) data.
\end{itemize}

\section*{Experience}
\noindent{\textbf{Center for Translational Research in Neuroimaging and Data Science}} \hfill January 2022 --- Now\\
Graduate Research Assistant \hfill \textit{Atlanta, USA}\\
[-2em]
\begin{itemize}[leftmargin=*,noitemsep]
    \item Built and evaluated machine learning models (classification, deep learning) on large-scale structured and unstructured datasets, achieving state-of-the-art performance.
    \item Performed data cleaning, preprocessing, and exploratory data analysis on large multimodal datasets to surface actionable insights.
    \item Developed reusable data pipelines and analysis tooling, automating ingestion, preprocessing, training, and evaluation.
    \item Collaborated with cross-functional teams (researchers and clinicians) to translate findings into practical tools and clear reports.
    \item Communicated technical findings to both technical and general audiences at internal meetings and international conferences.
\end{itemize}

\noindent\textbf{Institute of Radioengineering and Electronics} \hfill May 2017 --- September 2020\\
Research Assistant \hfill \textit{Moscow, Russia}\\
[-2em]
\begin{itemize}[leftmargin=*,noitemsep]
    \item Analyzed and modeled complex signal and physical systems using mathematical and numerical methods.
    \item Ran large-scale numerical simulations and analyzed results to validate theoretical models.
    \item Developed models providing theoretical validation for prototype hardware.
\end{itemize}

\section*{Education}
\noindent\textbf{Georgia State University} \hfill June 2027\\
PhD in Computer Science \hfill
% Machine Learning \& Neuroimaging,
GPA 3.89\\[0.3em]
\textbf{Moscow Institute of Physics and Technology} \hfill June 2020\\
MS in Applied Mathematics and Physics \hfill
% Spintronics,
GPA 3.92\\[0.3em]
\textbf{Moscow Institute of Physics and Technology} \hfill June 2018\\
BS in Applied Mathematics and Physics \hfill
% Solid State Physics,
GPA 3.81

\end{document}
